\documentclass{article}
\usepackage{lmodern}
\usepackage[T1]{fontenc}
\usepackage{amsmath,amssymb,amsthm,xcolor}
\newcommand{\verifierissue}[1]{%
  \par\noindent\textcolor{red}{\textbf{[Verifier issue:]}\ #1}\par\medskip}


\newtheorem{theorem}{Theorem}[section]
\newtheorem{lemma}[theorem]{Lemma}
\newtheorem{corollary}[theorem]{Corollary}
\newtheorem{proposition}[theorem]{Proposition}
\theoremstyle{definition}
\newtheorem{definition}[theorem]{Definition}
\newtheorem{remark}[theorem]{Remark}

\title{Spectral Lower Bound and an Improved Upper Bound\\ for the Chromatic Number of Unit Sphere Packings}
\date{}

\begin{document}
\maketitle

\begin{abstract}
A \emph{unit sphere packing} in $\mathbb{R}^n$ consists of unit spheres of radius $\tfrac12$ with pairwise disjoint interiors, equivalently a set of centres with pairwise distance $\ge 1$.  The \emph{tangency graph} connects pairs of centres at distance exactly $1$, and the \emph{chromatic number} of a packing is the chromatic number of its tangency graph.  Let $M(n)$ denote the supremum of chromatic numbers of unit sphere packings in $\mathbb{R}^n$.

\textbf{Lower bound.}  We give a self-contained spectral proof that
\[
  \limsup_{n\to\infty}\;\frac{M(n)-n}{n^{2/3}}\;\ge\;1,
\]
and more precisely $M(d_q)\ge q^3+1 = d_q+(q^2-q+1)$ for every prime power $q\ge 2$, where $d_q=q^3-q^2+q$.  The proof combines: (i) a \emph{Spectral Realization Theorem} showing that any connected $k$-regular graph with least adjacency eigenvalue $s<-1$ of multiplicity $g$ embeds as the tangency graph of a unit sphere packing in $\mathbb{R}^{v-g-1}$, and (ii) the classical \emph{Hoffman ratio bound} $\chi(G)\ge 1-k/s$.  The witness is the complement of the collinearity graph of $\mathrm{GQ}(q,q^2)$.

\textbf{Upper bound.}  We prove an improved upper bound
\[
  M(n)\;\le\; O\!\left(\tau(n)\cdot\sqrt{\frac{\log n}{n}}\right)
  \;=\;2^{0.401n(1+o(1))}\cdot O\!\left(\sqrt{\frac{\log n}{n}}\right),
\]
where $\tau(n)$ is the $n$-dimensional kissing number.  This improves the standard greedy bound $M(n)\le\tau(n)+1$ by a polynomial factor $O(\sqrt{n/\log n})\to\infty$.  The key inputs are: the geometric fact that the clique number of any packing's tangency graph satisfies $\omega\le n+1$ (equidistant-set bound), the degree bound $\Delta\le\tau(n)$, and the Bonamy--Kelly--Nelson--Postle theorem on coloring graphs with bounded maximum degree and small clique number.

\textbf{Barrier.}  A systematic barrier theorem (Section~6) proves that the four classical spectral families from finite geometry (Kneser graphs, triangular graphs, symplectic polar graphs, and complements of hexagon point graphs) all have Hoffman bound $H(G)\le d$, so none can improve the lower-bound surplus beyond $O(d^{2/3})$ within this framework.
\end{abstract}

\section*{Problem}

Throughout this paper, a \emph{unit sphere} has radius $\tfrac{1}{2}$ (diameter $1$).  Two such spheres have disjoint interiors if and only if their centres are at distance $\ge 1$; they are \emph{tangent} when their centres are at distance exactly $1$.

A \emph{unit sphere packing} in $\mathbb{R}^n$ is a set of unit spheres (radius $\tfrac{1}{2}$) with pairwise disjoint interiors, equivalently a set of points (sphere centres) in $\mathbb{R}^n$ with pairwise distance $\ge 1$. The \emph{tangency graph} of the packing takes the centres as vertices and connects pairs at distance exactly $1$ by an edge. A \emph{colouring} assigns a colour to each sphere so that tangent pairs receive different colours; the \emph{chromatic number} of the packing is the minimum number of colours needed.

Denote by $M(n)$ the supremum of chromatic numbers of all unit sphere packings (finite or infinite) in $\mathbb{R}^n$. The current state of knowledge is:
\begin{itemize}
  \item \textit{Upper bound}: $M(n)\le \tau(n)+1\le 2^{(0.401+o(1))n}$ by greedy colouring and the Kabatiansky--Levenshtein bound on the kissing number~\cite{KL1978}.
  \item \textit{Lower bound}: $M(n)\ge n+\Omega(n^{2/3})$ for an infinite sequence of dimensions $n$, first proved by Hao Chen~\cite{Chen2015} via generalized quadrangles.
\end{itemize}
We provide: (a) a clean spectral proof of the lower bound establishing the explicit constant $\limsup_{n\to\infty}(M(n)-n)/n^{2/3}\ge 1$; (b) an improved upper bound $M(n)=O(\tau(n)\sqrt{\log n/n})$, a polynomial-factor improvement over greedy; and (c) a barrier theorem for the lower-bound framework.

\section*{Main Results}

\subsection*{Section 1: Spectral Realization as a Unit Packing}

\begin{theorem}[Spectral Realization]
\label{thm:spectral-realization}
Let $G=(V,E)$ be a finite connected non-complete $k$-regular graph on $v$ vertices with adjacency matrix $A$, least eigenvalue $s<-1$ of multiplicity $g$, and $d=v-g-1$. Then there exist points $\{p_x : x\in V\}\subset \mathbb{R}^d$ such that
\[
  \|p_x-p_y\|=1\quad\text{if }xy\in E,
  \qquad
  \|p_x-p_y\|>1\quad\text{if }xy\notin E,\; x\neq y.
\]
In particular, the tangency graph of the unit sphere packing $\{p_x\}$ (with minimum inter-centre distance $1$) is exactly $G$.
\end{theorem}

We need to include one extra positive scale factor. The unscaled matrix $PMP$ gives the right \textbf{relative} distances, but not generally unit edge length unless $s=-2$. We define the scaled Gram matrix accordingly.

\begin{proof}
Let $\mathbf 1\in\mathbb R^v$ be the all-ones vector, let $J=\mathbf 1\mathbf 1^{\mathsf T}$, and let
\[
P=I-\frac{1}{v}J.
\]
Thus $P$ is the orthogonal projection onto $\mathbf 1^\perp$. Set
\[
M=\frac{1}{-s}(A-sI),
\qquad
\alpha=\frac{s}{2(s+1)}.
\]
Since $s<-1$, we have $\alpha>0$. We will use
\[
B=\alpha PMP
\]
as the Gram matrix.

Because $s$ is the least eigenvalue of the symmetric matrix $A$, every eigenvalue of $A-sI$ is nonnegative. Since $-s>0$, it follows that $M\succeq 0$. Hence for every $z\in\mathbb R^v$,
\[
z^{\mathsf T}PMPz=(Pz)^{\mathsf T}M(Pz)\ge 0,
\]
so $PMP\succeq 0$, and therefore $B\succeq 0$.

\textbf{Rank computation.} Since $G$ is $k$-regular, $A\mathbf 1=k\mathbf 1$. Moreover, because $G$ is connected, the $k$-eigenspace is exactly $\operatorname{span}\{\mathbf 1\}$. (Proof: if $Au=ku$ and $u_x:=u(x)$ achieves its maximum at some vertex $x_0$, then $ku_{x_0}=\sum_{y\sim x_0}u_y\le ku_{x_0}$ with equality only if all neighbours of $x_0$ also achieve the maximum; by connectedness all entries of $u$ equal $u_{x_0}$, so $u\in\operatorname{span}\{\mathbf 1\}$.)

Take an orthogonal eigenbasis of $A$. The vector $\mathbf 1$ contributes one eigenvector with eigenvalue $k$. The eigenspace for the least eigenvalue $s$ has dimension $g$ and lies inside $\mathbf 1^\perp$ (because for any eigenvector $u$ with $Au = s u$ and $s \ne k$, we have $k(1^T u) = (k\mathbf{1})^T u = (A\mathbf{1})^T u = \mathbf{1}^T A u = s(\mathbf{1}^T u)$, so $(k-s)(\mathbf{1}^T u)=0$, hence $\mathbf{1}^T u = 0$). On $\mathbf 1^\perp$, the projection $P$ acts as the identity, so if $u\perp \mathbf 1$ and $Au=\lambda u$, then
\[
PMPu=Mu=\frac{\lambda-s}{-s}u.
\]
This eigenvalue is $0$ exactly when $\lambda=s$, and is positive when $\lambda>s$. Also, $PMP\mathbf 1=0$. Therefore,
\[
\ker(PMP)=\operatorname{span}\{\mathbf 1\}\oplus \ker(A-sI),
\]
so $\dim\ker(PMP)=1+g$, and
\[
\operatorname{rank}(B)=\operatorname{rank}(PMP)=v-(1+g)=v-g-1=d.
\]

\textbf{Gram factorization.} Since $B$ is a real positive-semidefinite matrix of rank $d$, it admits a Gram factorisation: let $B=U\Delta U^{\mathsf T}$ be its eigendecomposition, where $\Delta=\mathrm{diag}(\beta_1,\ldots,\beta_d,0,\ldots,0)$ with $\beta_1,\ldots,\beta_d>0$ the positive eigenvalues, and $U$ orthogonal. Set $F=[{\sqrt{\beta_1}}\,u_1\mid \cdots\mid {\sqrt{\beta_d}}\,u_d]\in\mathbb R^{v\times d}$, so that $FF^{\mathsf T}=\sum_{i=1}^d\beta_i u_i u_i^{\mathsf T}=B$. For each vertex $x\in V$, define $p_x\in\mathbb R^d$ to be the $x$-th row of $F$. Then $B_{xy}=\langle p_x,p_y\rangle$.

\textbf{Distance computation.} Let
\[
\rho=\frac{k-s}{-s}.
\]
Every row sum of $M$ equals $\rho$: since $\sum_y A_{xy}=k$ (regularity) and $\sum_y I_{xy}=1$,
\[
\sum_y M_{xy}=\frac{1}{-s}(k-s)=\rho.
\]
Since $M$ is symmetric, every column sum also equals $\rho$. Therefore $M\mathbf{1}=\rho\mathbf{1}$, giving $MJ=\rho J$ and $JM=\rho J$, so $JMJ=J(\rho\mathbf{1}\mathbf{1}^T)=\rho v J$. Expanding:
\[
PMP=\Bigl(I-\tfrac{1}{v}J\Bigr)M\Bigl(I-\tfrac{1}{v}J\Bigr)
=M-\tfrac{1}{v}MJ-\tfrac{1}{v}JM+\tfrac{1}{v^2}JMJ
=M-\tfrac{\rho}{v}J-\tfrac{\rho}{v}J+\tfrac{\rho}{v}J
=M-\frac{\rho}{v}J.
\]
For every vertex $x$, $M_{xx}=\frac{-s}{-s}=1$. For distinct vertices $x\neq y$,
\[
M_{xy}=
\begin{cases}
\dfrac{1}{-s}, & xy\in E,\\[4pt]
0, & xy\notin E.
\end{cases}
\]
Therefore $(B_0)_{xx}=1-\frac{\rho}{v}$ (where $B_0=PMP$), and for $x\neq y$,
\[
(B_0)_{xy}=
\begin{cases}
-\dfrac{1}{s}-\dfrac{\rho}{v}, & xy\in E,\\[6pt]
-\dfrac{\rho}{v}, & xy\notin E.
\end{cases}
\]
Now let $x\neq y$. Since $B=\alpha B_0$ is the Gram matrix of the $p_x$'s,
\[
\|p_x-p_y\|^2 = B_{xx}+B_{yy}-2B_{xy}.
\]

If $xy\in E$:
\[
\|p_x-p_y\|^2
=\alpha\Bigl[2\Bigl(1-\tfrac{\rho}{v}\Bigr)-2\Bigl(-\tfrac{1}{s}-\tfrac{\rho}{v}\Bigr)\Bigr]
=\alpha\Bigl(2+\tfrac{2}{s}\Bigr)
=\frac{s}{2(s+1)}\cdot\frac{2(s+1)}{s}=1.
\]
Thus $\|p_x-p_y\|=1$ whenever $xy\in E$.

If $xy\notin E$, $x\neq y$:
\[
\|p_x-p_y\|^2
=\alpha\Bigl[2\Bigl(1-\tfrac{\rho}{v}\Bigr)-2\Bigl(-\tfrac{\rho}{v}\Bigr)\Bigr]
=2\alpha=\frac{s}{s+1}.
\]
Since $s<-1$, write $s=-1-\varepsilon$ with $\varepsilon>0$; then $\frac{s}{s+1}=\frac{1+\varepsilon}{\varepsilon}>1$. Hence $\|p_x-p_y\|>1$ whenever $xy\notin E$ and $x\neq y$.

Therefore the point set $\{p_x:x\in V\}\subset\mathbb R^d$ has minimum pairwise distance at least $1$, with equality exactly on edges of $G$. Consequently, the tangency graph of this unit sphere packing is exactly $G$.
\end{proof}

\subsection*{Section 2: Hoffman Lower Bound for Chromatic Number}

\begin{theorem}[Hoffman Bound]
\label{thm:hoffman}
Let $G$ be a finite $k$-regular graph on $v$ vertices with least adjacency eigenvalue $s<0$. Then
\[
  \chi(G) \;\ge\; 1-\frac{k}{s}.
\]
\end{theorem}

\begin{proof}
For reference, the inequality to be proved is the standard Hoffman lower bound on the chromatic number: $\chi(G)\ge 1-\lambda_{1}/\lambda_{n}$, where $\lambda_{1}$ and $\lambda_{n}$ are the largest and smallest adjacency eigenvalues. We prove the stated $k$-regular case directly.

Let $A$ be the adjacency matrix of $G$, and let $\mathbf{1}\in\mathbb{R}^{v}$ be the all-ones vector. Since $G$ is $k$-regular, $A\mathbf{1}=k\mathbf{1}$.

Let $S\subseteq V(G)$ be any independent set, and let $\mathbf{x}$ be its indicator vector. Put $n:=|S|=\mathbf{x}^{T}\mathbf{x}$. If $n=0$, there is nothing to prove, so assume $n>0$.

Decompose $\mathbf{x}$ into its component along $\mathbf{1}$ and its orthogonal component:
\[
\mathbf{x}=\frac{n}{v}\mathbf{1}+\mathbf{u},\qquad \mathbf{u}\perp \mathbf{1}.
\]
Indeed, $\mathbf{1}^{T}\mathbf{u}=\mathbf{1}^{T}\mathbf{x}-\frac{n}{v}\mathbf{1}^{T}\mathbf{1}=n-\frac{n}{v}v=0$.

Because $A$ is real symmetric, it has an orthonormal eigenbasis. Every eigenvalue of $A$ is at least $s$, so for every vector $\mathbf{u}\perp \mathbf{1}$ we have $\mathbf{u}^{T}A\mathbf{u}\ge s\|\mathbf{u}\|^{2}$.

Now compute $\mathbf{x}^{T}A\mathbf{x}$. Since $S$ is independent, no edge has both endpoints in $S$, so $\mathbf{x}^{T}A\mathbf{x}=0$. On the other hand,
\[
\mathbf{x}^{T}A\mathbf{x}
=\left(\frac{n}{v}\mathbf{1}+\mathbf{u}\right)^{T}A\left(\frac{n}{v}\mathbf{1}+\mathbf{u}\right)
=\frac{n^{2}}{v^{2}}\mathbf{1}^{T}A\mathbf{1}+\frac{2n}{v}\mathbf{1}^{T}A\mathbf{u}+\mathbf{u}^{T}A\mathbf{u}.
\]
Since $A\mathbf{1}=k\mathbf{1}$ and $\mathbf{u}\perp\mathbf{1}$, we have $\mathbf{1}^{T}A\mathbf{1}=kv$ and $\mathbf{1}^{T}A\mathbf{u}=(A\mathbf{1})^{T}\mathbf{u}=k\mathbf{1}^{T}\mathbf{u}=0$.
Therefore
\[
0=\mathbf{x}^{T}A\mathbf{x}=\frac{kn^{2}}{v}+\mathbf{u}^{T}A\mathbf{u}.
\]
Using $\mathbf{u}^{T}A\mathbf{u}\ge s\|\mathbf{u}\|^{2}$ gives
\[
0\ge\frac{kn^{2}}{v}+s\|\mathbf{u}\|^{2}.
\]
Now $\|\mathbf{u}\|^{2}=\|\mathbf{x}-\frac{n}{v}\mathbf{1}\|^{2}=\|\mathbf{x}\|^{2}-\frac{n^{2}}{v}=n-\frac{n^{2}}{v}=\frac{n(v-n)}{v}$.
Hence
\[
0\ge\frac{kn^{2}}{v}+s\frac{n(v-n)}{v}.
\]
Multiplying by $\frac{v}{n}>0$ yields
\[
0\ge kn+s(v-n)=(k-s)n+sv.
\]
Thus $(k-s)n\le -sv$. Since $s<0$ and $k-s>0$, we obtain
\[
n\le \frac{-sv}{k-s}.
\]
So every independent set in $G$ has size at most $\frac{-sv}{k-s}$.

Now let $V_{1},\dots,V_{\chi}$ be the color classes of an optimal proper coloring of $G$, where $\chi=\chi(G)$. Each $V_{i}$ is an independent set, so $|V_{i}|\le \frac{-sv}{k-s}$ for each $i$. Since the color classes partition $V(G)$,
\[
v=\sum_{i=1}^{\chi}|V_{i}|\le\chi\cdot\frac{-sv}{k-s}.
\]
Dividing by $v>0$ gives $1\le \chi\cdot\frac{-s}{k-s}$. Therefore
\[
\chi\ge\frac{k-s}{-s}=1-\frac{k}{s}.\qquad\square
\]
\end{proof}

\subsection*{Section 3: The Generalized Quadrangle Construction}

\begin{theorem}[Spectral Family from $\mathrm{GQ}(q,q^2)$]
\label{thm:gq-family}
For every prime power $q\ge 2$, let $G_q$ denote the complement of the point-collinearity graph of the generalized quadrangle $\mathrm{GQ}(q,q^2)$. Then:
\begin{enumerate}
  \item $G_q$ has $v_q=(q+1)(q^3+1)$ vertices and is regular of degree $k_q=q^4$.
  \item The least adjacency eigenvalue of $G_q$ is $s_q=-q$, with multiplicity $g_q=q^4+q^2$, and the embedding dimension is $d_q:=v_q-g_q-1=q^3-q^2+q$.
  \item The Hoffman lower bound satisfies
    \[1-\frac{k_q}{s_q}=q^3+1=d_q+(q^2-q+1).\]
  \item For every prime power $q\ge 2$,
    \[1-\frac{k_q}{s_q}\ge d_q+\left(\frac{3}{4}\right)^{1/3}d_q^{2/3}.\]
\end{enumerate}
\end{theorem}

\begin{proof}
Let $\mathcal{Q}$ be a generalized quadrangle of order $(q,q^2)$, and let $\Gamma_q$ be its point-collinearity graph. Thus two distinct points are adjacent in $\Gamma_q$ exactly when they are collinear in $\mathcal{Q}$. Generalized quadrangles of order $(q,q^2)$ exist for all prime powers $q$: the classical realization is the \emph{elliptic quadric} $Q^-(5,q)$ in $\mathrm{PG}(5,q)$, which has $(q^3+1)(q+1)$ points and forms a $\mathrm{GQ}(q,q^2)$ (see \cite{PT1984} and \cite{BvM2022}). The Hermitian variety $H(3,q^2)$ gives $\mathrm{GQ}(q^2,q)$, not $\mathrm{GQ}(q,q^2)$.

The collinearity graph of a $\mathrm{GQ}(s,t)$ is strongly regular with parameters \cite{PT1984}
\[
v=(s+1)(st+1),\quad k=s(t+1),\quad \lambda=s-1,\quad \mu=t+1,
\]
and nontrivial eigenvalues $s-1$ (multiplicity $st(s+1)(t+1)/(s+t)$) and $-(t+1)$ (multiplicity $s^2(st+1)/(s+t)$).

Setting $s=q$ and $t=q^2$:
\[
v=(q+1)(q^3+1),\quad \kappa=q(q^2+1),\quad \lambda=q-1,\quad \mu=q^2+1,
\]
with nontrivial eigenvalues
\[
r=q-1,\qquad s_0=-(q^2+1).
\]

\textbf{Connectedness.} If two points are collinear, they are adjacent in $\Gamma_q$. If two points are not collinear, they have exactly $\mu=q^2+1>0$ common neighbours. Hence every pair of vertices is at distance at most $2$ in $\Gamma_q$, so $\Gamma_q$ is connected. For any connected $k$-regular graph, the degree eigenvalue $k$ has multiplicity $1$: if $Au=\kappa u$ then the maximum-entry argument (same as in Section~1) shows $u\in\operatorname{span}\{\mathbf{1}\}$, so the $\kappa$-eigenspace is 1-dimensional.

\textbf{Multiplicity computation.} Let the multiplicities of $r$ and $s_0$ be $f$ and $g_0$ respectively. The trace of the adjacency matrix of a simple graph (no self-loops) equals $0$ since $\operatorname{tr}(A)=\sum_x A_{xx}=0$; by the spectral theorem, this equals the sum of eigenvalues with multiplicity: $\kappa\cdot 1 + r\cdot f + s_0\cdot g_0 = 0$. Combined with $1+f+g_0=v$, the trace condition is $\kappa+(q-1)f-(q^2+1)g_0=0$. Using $g_0=v-1-f$:
\[
\kappa+(q-1)f-(q^2+1)(v-1-f)=0 \implies (q^2+q)f=(q^2+1)(v-1)-\kappa.
\]
Substituting $v-1=q^4+q^3+q$ and $\kappa=q(q^2+1)$:
\[
f=\frac{(q^2+1)(q^4+q^3+q)-q(q^2+1)}{q^2+q}
=\frac{(q^2+1)(q^4+q^3)}{q(q+1)}
=q^2(q^2+1)=q^4+q^2.
\]
Hence $g_0=v-1-f=(q^4+q^3+q)-(q^4+q^2)=q^3-q^2+q$.

\textbf{Complement eigenvalues.} Let $G_q=\overline{\Gamma_q}$ with adjacency matrix $B=J-I-A$. Since $\Gamma_q$ is $\kappa$-regular, $B\mathbf{1}=(v-1-\kappa)\mathbf{1}$, so $G_q$ has degree
\[
k_q=v-1-\kappa=(q^4+q^3+q)-(q^3+q)=q^4.
\]
This proves item (1).

For the non-principal eigenvalues: if $Au=\theta u$ with $\theta\ne\kappa$, then $\kappa(\mathbf{1}^Tu)=\mathbf{1}^T A u=\theta(\mathbf{1}^Tu)$, so $\mathbf{1}^Tu=0$. Thus all non-principal eigenvectors lie in $\mathbf{1}^\perp$, and for $u\perp\mathbf{1}$ with $Au=\theta u$: $Bu=(J-I-A)u=0-u-\theta u=(-1-\theta)u$. Hence $G_q$ has eigenvalues:
\begin{itemize}
  \item $q^4$ (multiplicity $1$, principal eigenvalue),
  \item $-r-1=-q$ (multiplicity $f=q^4+q^2$),
  \item $-s_0-1=q^2$ (multiplicity $g_0=q^3-q^2+q$).
\end{itemize}

For $q\ge 2$, the three eigenvalues satisfy $q^4>0$, $q^2>0$, and $-q<0$. Hence the least adjacency eigenvalue of $G_q$ is $s_q=-q$ with multiplicity $g_q=f=q^4+q^2$.

\textbf{Embedding dimension.}
\[
d_q:=v_q-g_q-1=(q^4+q^3+q+1)-(q^4+q^2)-1=q^3-q^2+q.
\]

\textbf{Non-completeness and connectivity of $G_q$.} We have $k_q=q^4<q^4+q^3+q=v_q-1$, so $G_q$ is not complete.

For connectivity, recall that two vertices $x,y$ are \emph{adjacent in $G_q$} if and only if they are \emph{non-collinear in $\mathcal{Q}$}. We show every pair $x\ne y$ has a common neighbour in $G_q$, establishing diameter at most~$2$.

\emph{Case 1: $x$ and $y$ are adjacent in $G_q$} (non-collinear in $\mathcal{Q}$). Then $xy$ is already an edge, and they trivially have a common neighbour if we can find $z\ne x,y$ adjacent to both; the diameter-2 bound suffices for connectivity, but adjacency already gives a path.

\emph{Case 2: $x$ and $y$ are non-adjacent in $G_q$} (collinear in $\mathcal{Q}$). We need a vertex $z$ adjacent to both $x$ and $y$ in $G_q$, i.e., non-collinear with both. In $\mathcal{Q}$, each point has $\kappa=q(q^2+1)$ collinear points, so $x$ is non-collinear with exactly $v-1-\kappa=q^4$ points of $\mathcal{Q}\setminus\{x\}$; call this set $A$. Similarly let $B$ be the $q^4$ points non-collinear with $y$. Since $x$ and $y$ are collinear, $x\notin A$ and $y\notin B$ (a point is always collinear with itself, but the exclusion $\{x,y\}$ merely removes two points from the ambient set). More precisely, $A\subseteq V\setminus\{x\}$ and $B\subseteq V\setminus\{y\}$, and since $x$ is collinear with $y$ we have $x\notin B$ and $y\notin A$. So both $A$ and $B$ are subsets of $V\setminus\{x,y\}$, which has $v_q-2=q^4+q^3+q-1$ elements.

Since $|A|+|B|=2q^4>q^4+q^3+q-1=|V\setminus\{x,y\}|$ for all $q\ge 2$ (because $q^4-(q^3+q-1)=(q-1)(q^3-1)>0$ for $q\ge 2$), by inclusion-exclusion $A\cap B\ne\emptyset$. Any $z\in A\cap B$ is non-collinear with both $x$ and $y$, so $z$ is adjacent to both in $G_q$, giving the path $x$-$z$-$y$.

Thus $G_q$ has diameter at most $2$ and is connected. This proves item (2).

\textbf{Hoffman bound.}
\[
1-\frac{k_q}{s_q}=1-\frac{q^4}{-q}=1+q^3=q^3+1.
\]
Since $d_q+(q^2-q+1)=(q^3-q^2+q)+(q^2-q+1)=q^3+1$, we have
\[
1-\frac{k_q}{s_q}=d_q+(q^2-q+1).
\]
This proves item (3).

\textbf{Explicit surplus bound.} Let $A_q=q^2-q+1$, so $d_q=qA_q$ and the surplus equals $A_q$. It suffices to show $A_q\ge(3/4)^{1/3}d_q^{2/3}=(3/4)^{1/3}(qA_q)^{2/3}$. Since $A_q>0$, cubing gives the equivalent inequality $A_q^3\ge(3/4)q^2 A_q^2$, i.e.\ $A_q\ge(3/4)q^2$. Computing:
\[
A_q-\frac{3}{4}q^2=q^2-q+1-\frac{3}{4}q^2=\frac{1}{4}q^2-q+1=\frac{(q-2)^2}{4}\ge 0.
\]
Therefore $q^2-q+1\ge(3/4)^{1/3}d_q^{2/3}$ for all $q$, and item (4) follows.
\end{proof}

\subsection*{Section 4: Packing Lower Bounds at the Special Dimensions}

\begin{theorem}[Main Packing Bound]
\label{thm:main-packing}
For every prime power $q\ge 2$, the supremum $M(d_q)$ of chromatic numbers of finite unit sphere packings in $\mathbb{R}^{d_q}$ satisfies
\[
  M(d_q)\;\ge\; q^3+1 \;=\; d_q + (q^2-q+1) \;\ge\; d_q+\left(\frac{3}{4}\right)^{1/3}d_q^{2/3},
\]
where $d_q = q^3-q^2+q$.
\end{theorem}

\begin{proof}
Fix a prime power $q\ge 2$. Let $G_q$ be the graph from Section~3. We explicitly verify that $G_q$ meets all the hypotheses of the Spectral Realization Theorem (Section~1):
\begin{itemize}
  \item \emph{Finite}: $G_q$ has $v_q=(q+1)(q^3+1)$ vertices (Section~3, item~1).
  \item \emph{Connected}: proved in Section~3 (diameter-at-most-2 argument via the $|A\cap B|\ne\emptyset$ counting for non-adjacent pairs).
  \item \emph{Non-complete}: $k_q=q^4<v_q-1$ (Section~3, item~1).
  \item \emph{$k_q$-regular}: $G_q$ is regular of degree $k_q=q^4$ (Section~3, item~1).
  \item \emph{Least eigenvalue $s_q=-q<-1$}: since $q\ge 2$, we have $s_q=-q\le -2<-1$ (Section~3, item~2).
  \item \emph{Multiplicity $g_q$}: $g_q=q^4+q^2$ (Section~3, item~2).
  \item \emph{Dimension formula}: $d_q=v_q-g_q-1=q^3-q^2+q$ (Section~3, item~2).
\end{itemize}
By the Spectral Realization Theorem, there exist points $\{p_x:x\in V(G_q)\}\subset\mathbb{R}^{d_q}$ forming a unit sphere packing in $\mathbb{R}^{d_q}$ whose tangency graph is exactly $G_q$. Therefore
\[
M(d_q)\ge\chi(G_q).
\]

Since $G_q$ is $k_q$-regular with least eigenvalue $s_q=-q<0$, the Hoffman Bound (Section~2) gives
\[
\chi(G_q)\ge 1-\frac{k_q}{s_q}=1+q^3=q^3+1.
\]

By Section~3, item~3 and item~4:
\[
q^3+1=d_q+(q^2-q+1)\ge d_q+\left(\frac{3}{4}\right)^{1/3}d_q^{2/3}.
\]

Combining: $M(d_q)\ge\chi(G_q)\ge q^3+1\ge d_q+(3/4)^{1/3}d_q^{2/3}$.
\end{proof}

\subsection*{Section 5: Monotonicity and the Limsup Lower Bound}

\begin{theorem}[Monotonicity and Limsup Bound]
\label{thm:limsup}
Let $M(n)$ denote the supremum of chromatic numbers of unit sphere packings (finite or infinite) in $\mathbb{R}^n$. Then:
\begin{enumerate}
  \item \textbf{(Monotonicity)} $M(n+1)\ge M(n)$ for all $n\ge 1$.
  \item \textbf{(Limsup lower bound)}
    \[
      \limsup_{n\to\infty}\frac{M(n)-n}{n^{2/3}}\ge 1.
    \]
\end{enumerate}
\end{theorem}

\begin{remark}
The constant $1$ is approached but not necessarily achieved: the proof gives $(M(d_q)-d_q)/d_q^{2/3} \ge (1-1/q+1/q^2)^{1/3}$, which tends to $1$ from below as $q\to\infty$.  No upper bound of the form $M(n)-n = O(n^{2/3})$ is claimed here.
\end{remark}

\begin{remark}
The gap $d_{q+1}-d_q=3q^2+q+1\sim 3d_q^{2/3}$ is of the same order as the surplus $q^2-q+1\sim d_q^{2/3}$, so a naive monotonicity interpolation does \emph{not} extend the bound $M(n)\ge n+c\,n^{2/3}$ to all large $n$ from this sequence alone.
\end{remark}

\begin{proof}
\textbf{(1) Monotonicity.}
Define the embedding $\iota:\mathbb{R}^n\to\mathbb{R}^{n+1}$ by $\iota(p)=(p,0)$.  For any unit sphere packing $\mathcal{P}$ in $\mathbb{R}^n$ (finite or infinite), the image $\iota(\mathcal{P})=\{(p,0):p\in\mathcal{P}\}$ is a unit sphere packing in $\mathbb{R}^{n+1}$, because
\[
\|\iota(p)-\iota(q)\|_{\mathbb{R}^{n+1}}=\|p-q\|_{\mathbb{R}^n}
\]
for all $p,q\in\mathcal{P}$.  In particular, distances are preserved, so tangencies are preserved, and the tangency graph of $\iota(\mathcal{P})$ is isomorphic to that of $\mathcal{P}$.  Hence every chromatic number achievable by a packing in $\mathbb{R}^n$ is also achievable by a packing in $\mathbb{R}^{n+1}$, giving $M(n)\le M(n+1)$.

\textbf{(2) Limsup bound.}
By the Main Packing Bound (Section~4), for every prime power $q\ge 2$:
\[
M(d_q)\;\ge\; q^3+1 \;=\; d_q + (q^2-q+1),\qquad d_q=q^3-q^2+q.
\]
Since $d_q=q(q^2-q+1)$, we have $d_q^{2/3}=q^{2/3}(q^2-q+1)^{2/3}$, so
\[
\frac{M(d_q)-d_q}{d_q^{2/3}}\;\ge\;\frac{q^2-q+1}{d_q^{2/3}}=\frac{(q^2-q+1)^{1/3}}{q^{2/3}}=\left(1-\frac{1}{q}+\frac{1}{q^2}\right)^{1/3}.
\]
The function $f(q)=1-1/q+1/q^2$ satisfies $f'(q)=1/q^2-2/q^3=(q-2)/q^3\ge 0$ for $q\ge 2$, so $f$ is nondecreasing for $q\ge 2$ and $\lim_{q\to\infty}f(q)=1$.

Fix any $\varepsilon\in(0,1)$.  Since $f(q)\to 1$, there exists a prime power $q_0$ such that $f(q_0)>1-\varepsilon$.  By Euclid's theorem on the infinitude of primes, there exist infinitely many prime powers $q>q_0$, giving infinitely many indices $n_j=d_{q_j}\to\infty$ with
\[
\frac{M(n_j)-n_j}{n_j^{2/3}}\;\ge\;f(q_j)^{1/3}\;>\;(1-\varepsilon)^{1/3}.
\]
Therefore $\limsup_{n\to\infty}(M(n)-n)/n^{2/3}\ge (1-\varepsilon)^{1/3}$ for every $\varepsilon\in(0,1)$.  Letting $\varepsilon\to 0$ gives
\[
\limsup_{n\to\infty}\frac{M(n)-n}{n^{2/3}}\;\ge\;1.\qquad\square
\]
\end{proof}

\subsection*{Section 6: Barrier Theorem for the Hoffman Framework}

The main result gives $M(n)\ge n+\Omega(n^{2/3})$ via the Spectral Realization Theorem and the Hoffman bound.  A natural question is whether the same framework can yield a larger surplus of order $n^\alpha$ for some $\alpha>2/3$.  This section shows that for all standard graph families from finite geometry the answer is no.

\medskip
We collect the framework parameters.  For a connected $k$-regular graph $G$ with $v$ vertices and least eigenvalue $s<-1$ of multiplicity $g$:
\begin{itemize}
  \item embedding dimension: $d = v - g - 1$,
  \item Hoffman lower bound: $H(G) := 1 - k/s$,
  \item surplus: $H(G) - d$.
\end{itemize}

\begin{remark}
Throughout this section, ``standard spectral families from finite geometry'' refers to exactly the four families listed below: Kneser graphs, triangular graphs, non-collinearity graphs of symplectic polar spaces, and complements of generalized-hexagon point graphs.  The theorem makes no claim about other families.
\end{remark}

\begin{theorem}[Hoffman-Framework Barrier]\label{thm:barrier}
For the following four families, the Hoffman bound does not exceed $d$:
\begin{enumerate}
  \item \textup{(Kneser $K(n,k)$, fixed $k\ge 3$):} $H = n/k$, $d\sim n^k/k!$, so $H = o(d)$.
  \item \textup{(Triangular $T(n)$, $n\ge 5$):} $H = d = n-1$, surplus $= 0$.
  \item \textup{(Non-collinearity of $W(2r{-}1,q)$, $r\ge 2$):} $H\le d$, equality only at $(r,q)=(2,2)$.
  \item \textup{(Hexagon complement, order $(q,q)$):} $H\sim\tfrac12q^4$, $d\sim\tfrac56q^5$, so $H = o(d)$.
\end{enumerate}
None of these families achieves $H(G)\ge d+c\,d^\alpha$ for any $c>0$, $\alpha>0$.
\end{theorem}

\begin{proof}
\textbf{(1) Kneser graph $K(n,k)$, fixed $k\ge 3$, $n > 2k$.}
The Kneser graph $K(n,k)$ has vertex set $\binom{[n]}{k}$ (all $k$-subsets of $[n]$), adjacent when disjoint.  For $n > 2k$, its adjacency eigenvalues are
\[
  \lambda_j = (-1)^j\binom{n-k-j}{k-j},\quad 0\le j\le k,
\]
with multiplicities $\binom{n}{j}-\binom{n}{j-1}$ for $j\ge 1$~\cite{BvM2022}.  In particular the $j=1$ term,
\[
  \lambda_1 = -\binom{n-k-1}{k-1},
\]
has multiplicity $\binom{n}{1}-\binom{n}{0} = n-1$.

\textit{Claim}: $\lambda_1$ is the least eigenvalue for every $n > 2k$, $k\ge 3$.

\textit{Proof}: The even-$j$ eigenvalues are positive for $n>2k$.  For odd $j\ge 3$, compare with $j=1$:
\[
  \frac{|\lambda_1|}{|\lambda_j|}
  = \frac{\binom{n-k-1}{k-1}}{\binom{n-k-j}{k-j}}
  = \prod_{i=1}^{j-1}\frac{n-k-i}{k-i}.
\]
For each $i\in\{1,\ldots,j-1\}$: since $j\le k$ implies $i\le k-2$, the denominator $k-i\ge 2 > 0$.  Since $n>2k$, we have $n-k-i > k-i$ (equivalently $n>2k$), so each factor exceeds $1$.  Hence the product exceeds $1$, giving $|\lambda_j| < |\lambda_1|$ for every odd $j\ge 3$.  Thus $s_G = \lambda_1$ with $g = n-1$.\hfill$\square$

The degree is $k_G = \binom{n-k}{k}$, so the Hoffman bound
$H = 1 + \binom{n-k}{k}/\binom{n-k-1}{k-1} = n/k$,
and $d = \binom{n}{k}-n\sim n^k/k!$.  For fixed $k\ge 3$: $H/d \to 0$.

\textbf{(2) Triangular graph $T(n)$, $n\ge 5$.}
The triangular graph $T(n)$ has vertex set $\binom{[n]}{2}$ with two vertices adjacent when they share one element; $v=\binom{n}{2}$, $k_G=2(n-2)$.  Two distinct $2$-subsets of $[n]$ either are disjoint or share one element, so $T(n) = \overline{K(n,2)}$.

We derive the spectrum of $T(n)$ from the formula of Part (1) applied to $K(n,2)$ (with $k=2$~\cite{BvM2022}):
\begin{itemize}
  \item $j=1$: $\lambda_1 = -\binom{n-3}{1} = -(n-3)$, multiplicity $n-1$.
  \item $j=2$: $\lambda_2 = \binom{n-4}{0} = 1$, multiplicity $\binom{n}{2}-n$.
\end{itemize}
Each nontrivial eigenvector of $K(n,2)$ lies in $\mathbf{1}^\perp$ and becomes an eigenvector of $T(n)=\overline{K(n,2)}$ with eigenvalue $-1-\lambda$, giving:
\[
  s_G = -1-1 = -2\;\text{ (mult }g=\tbinom{n}{2}-n\text{)},
  \qquad
  n-4\;\text{ (mult }n-1\text{)}.
\]
Hence $d = \binom{n}{2}-(\binom{n}{2}-n)-1 = n-1$ and $H = 1-2(n-2)/(-2) = n-1 = d$; surplus $= 0$.

\textbf{(3) Non-collinearity graph of $W(2r-1,q)$, $r\ge 2$.}
The symplectic polar space $W(2r-1,q)$ has $v=(q^{2r}-1)/(q-1)$ points.  Its non-collinearity graph $G$ is strongly regular with degree $k_G=q^{2r-1}$ and nontrivial eigenvalues $\pm q^{r-1}$~\cite{BvM2022}.  From the trace equations $f+g=v-1$ and $f-g=-k_G/q^{r-1}=-q^r$:
\[
  g = \tfrac{v-1+q^r}{2},\quad d = v-g-1 = \tfrac{v-1-q^r}{2},\quad
  H = 1 + q^r.
\]
We show $d\ge H$, i.e.\ $v\ge 3+3q^r$, by three cases.

\textit{Case A: $r=2$.}
$v = q^3+q^2+q+1$, so $v-3-3q^2 = q^3-2q^2+q-2 = (q-2)(q^2+1)\ge 0$,
with equality iff $q=2$.  At $(r,q)=(2,2)$: $d=H=5$.

\textit{Case B: $r=3$, $q=2$.}
Direct computation: $v=(2^6-1)/(2-1)=63$, $d=(63-1-8)/2=27$, $H=1+8=9$; $d=27 > H=9$.

\textit{Case C: $r\ge 3$, $(r,q)\ne(3,2)$.}
Here $q^{r-1}>4$: if $r=3$ then $q\ge 3$ gives $q^2\ge 9$; if $r\ge 4$ then $q^{r-1}\ge 2^3=8$.  Using $v\ge q^{2r-1}$ and $3+3q^r\le 4q^r$:
\[
  v-3-3q^r \ge q^{2r-1}-4q^r = q^r(q^{r-1}-4) > 0.
\]
Thus $d > H$ in all cases except $(r,q)=(2,2)$ where $d=H$.

\textbf{(4) Complement of the generalized hexagon point graph, order $(q,q)$.}
Let $\Gamma$ be the point-collinearity graph of a generalized hexagon of order $(q,q)$~\cite{BvM2022}: $v=(q+1)(q^4+q^2+1)$, degree $q(q+1)$, nontrivial eigenvalues $2q-1$ (mult $m_1$), $-1$ (mult $m_2$), $-(q+1)$ (mult $m_3$), where
\[
  m_1 = \tfrac{q(q+1)^2(q^2+q+1)}{6},\quad
  m_2 = \tfrac{q(q+1)^2(q^2-q+1)}{2},\quad
  m_3 = \tfrac{q(q^2+q+1)(q^2-q+1)}{3}.
\]
For $G=\overline{\Gamma}$: $k_G = q^3(q^2+q+1)$, and the nontrivial eigenvalues of $G$ are $-2q$ (mult $m_1$), $0$ (mult $m_2$), $q$ (mult $m_3$).  The least eigenvalue is $s_G=-2q$ with multiplicity $g=m_1$, so
\[
  d = v - m_1 - 1 = m_2 + m_3 \sim \tfrac56q^5,
  \qquad
  H = 1 + \tfrac{q^2(q^2+q+1)}{2} \sim \tfrac12q^4.
\]
Thus $H/d \sim 3/(5q)\to 0$.  Explicit checks: for $q=2$, $(m_1,m_2,m_3)=(21,27,14)$, $v=63$, $d=41$, $H=15$; for $q=3$, $d=259$, $H=59.5$.  In all cases $H < d$.

\medskip
\textbf{Summary.}  For Kneser ($k\ge 3$) and the hexagon complement, $H=o(d)$; for the triangular graph and symplectic cases, $H\le d$ with equality only at $(r,q)=(2,2)$.  None achieves $H\ge d+c\,d^\alpha$ for any $c,\alpha>0$.
\end{proof}

\begin{remark}
Any improvement to the surplus exponent $\alpha=2/3$ within the Hoffman framework requires either a genuinely new graph family or a proof technique beyond the Hoffman ratio bound.
\end{remark}

\subsection*{Section 7: Improved Upper Bound via Small-Clique Graph Coloring}

The following theorem improves the standard upper bound $M(n)\le \tau(n)+1$ by a polynomial factor, where $\tau(n)$ denotes the $n$-dimensional kissing number.

\begin{theorem}[Improved Upper Bound]\label{thm:upper-bound}
There exists an absolute constant $C>0$ such that, for all sufficiently large $n$,
\[
  M(n) \;\le\; C\,\tau(n)\,\sqrt{\frac{\ln(n+2)}{\ln \tau(n)}}.
\]
In particular, using the Kabatiansky--Levenshtein bound $\tau(n)\le 2^{0.401n(1+o(1))}$,
\[
  M(n) \;=\; O\!\left(\tau(n)\cdot\sqrt{\frac{\log n}{n}}\right)
  \;=\; 2^{0.401n(1+o(1))}\cdot O\!\left(\sqrt{\frac{\log n}{n}}\right).
\]
This is a polynomial-factor improvement over the greedy bound $M(n)\le\tau(n)+1$.
\end{theorem}

The proof rests on three ingredients, which we state as lemmas.

\begin{lemma}[de Bruijn--Erd\H{o}s Compactness]\label{lem:compactness}
$M(n) = \sup\bigl\{\chi(G_X) : X\subset\mathbb{R}^n\text{ finite},\;\|x-y\|\ge 1\ \forall x\ne y\in X\bigr\}.$
In other words, $M(n)$ is achieved as a supremum over \emph{finite} packings.
\end{lemma}

\begin{proof}
Let $\mathcal{P}$ be any unit sphere packing in $\mathbb{R}^n$ (possibly infinite), and let $G=G_{\mathcal{P}}$ be its tangency graph.  Every finite subgraph of $G$ is the tangency graph of a finite sub-packing, hence has chromatic number at most $\sup_{\text{finite }X}\chi(G_X)=:K$.  By the de Bruijn--Erd\H{o}s theorem~\cite{dBE1951}, an infinite graph whose every finite subgraph is $K$-colorable is itself $K$-colorable, so $\chi(G)\le K$.  Hence $M(n)\le K$; the reverse inequality $M(n)\ge K$ is trivial.
\end{proof}

\begin{lemma}[Geometric Graph Bounds]\label{lem:graph-bounds}
For every finite unit sphere packing $X\subset\mathbb{R}^n$ (centres with pairwise distance $\ge 1$), the tangency graph $G_X$ satisfies:
\begin{enumerate}
  \item[\textup{(a)}] \textup{(Degree bound)} $\Delta(G_X)\le\tau(n)$, the $n$-dimensional kissing number.
  \item[\textup{(b)}] \textup{(Clique bound)} $\omega(G_X)\le n+1$.
\end{enumerate}
\end{lemma}

\begin{proof}
\textbf{(a).} The neighbours of a centre $v\in X$ in $G_X$ are the centres $u\in X$ with $\|u-v\|=1$, all lying on the unit sphere $S^{n-1}(v)$ of radius $1$ centred at $v$, and pairwise at distance $\ge 1$.  Translating to $v=0$, this is a set of unit-sphere points with pairwise angular separation $\ge 60°$, i.e.\ a kissing configuration.  By definition of $\tau(n)$, the size of any kissing configuration is at most $\tau(n)$.

\textbf{(b).} A $k$-clique in $G_X$ is a set $\{x_1,\ldots,x_k\}\subset X$ with $\|x_i-x_j\|=1$ for all $i\ne j$ — an equidistant set of diameter $1$ in $\mathbb{R}^n$.  Setting $y_i=x_i-x_1$ for $i\ge 2$, we have $\|y_i\|=1$ and
\[
  \|y_i-y_j\|^2 = \|x_i-x_j\|^2 = 1, \quad i\ne j\ge 2,
\]
which gives $\langle y_i,y_j\rangle = \tfrac{1}{2}$ for all $i\ne j\ge 2$.  The Gram matrix $M$ of $y_2,\ldots,y_k$ satisfies $M_{ii}=1$ and $M_{ij}=\tfrac{1}{2}$ for $i\ne j$, i.e.\ $M=\tfrac{1}{2}I+\tfrac{1}{2}J$ where $J$ is the $(k-1)\times(k-1)$ all-ones matrix.  The eigenvalues of $M$ are $\tfrac{k}{2}$ (once) and $\tfrac{1}{2}$ ($k-2$ times), all positive, so $M$ is positive definite and has rank $k-1$.  Since $y_2,\ldots,y_k\in\mathbb{R}^n$, we need $k-1\le n$, i.e.\ $k\le n+1$.
\end{proof}

\begin{lemma}[Bonamy--Kelly--Nelson--Postle, 2022~{\cite{BKNP2022}}]\label{lem:bknp}
There exists an absolute constant $C_0>0$ such that the following holds.  For all $\omega\ge 3$ and all $\Delta\ge \omega^{C_0}$, every graph $G$ with $\Delta(G)\le\Delta$ and $\omega(G)\le\omega$ satisfies
\[
  \chi(G)\;\le\; 72\,\Delta\,\sqrt{\frac{\ln\omega}{\ln\Delta}}.
\]
\end{lemma}

\begin{proof}[Proof of Theorem~\ref{thm:upper-bound}]
By Lemma~\ref{lem:compactness}, $M(n)=\sup_{\text{finite }X}\chi(G_X)$.  Fix a finite unit sphere packing $X\subset\mathbb{R}^n$.  By Lemma~\ref{lem:graph-bounds}, $G_X$ has $\Delta(G_X)\le\tau(n)$ and $\omega(G_X)\le n+1$.

We verify the hypothesis of Lemma~\ref{lem:bknp} with $\Delta=\tau(n)$ and $\omega=n+2$: since $\tau(n)\ge 2^{0.2075n(1+o(1))}$ by the Kabatiansky--Levenshtein lower bound~\cite{KL1978}, we have $\ln\tau(n)\ge 0.2075n\ln 2$, while $\ln(n+2)^{C_0}=O(n)$, so $\tau(n)\ge(n+2)^{C_0}$ for all sufficiently large $n$.  Hence
\[
  \chi(G_X) \;\le\; 72\,\tau(n)\,\sqrt{\frac{\ln(n+2)}{\ln\tau(n)}}.
\]
Taking the supremum over all finite $X$ gives the first bound.

For the asymptotic form, since $\ln\tau(n)\le 0.401n\ln 2+o(n)$ (upper KL bound~\cite{KL1978}),
\[
  \sqrt{\frac{\ln(n+2)}{\ln\tau(n)}} = O\!\left(\sqrt{\frac{\log n}{n}}\right),
\]
giving $M(n)=O(\tau(n)\sqrt{\log n/n})=2^{0.401n(1+o(1))}\cdot O(\sqrt{\log n/n})$.

To see this is an improvement over greedy: the ratio of the new bound to $\tau(n)$ is $O(\sqrt{\log n/n})\to 0$, so the new bound is $o(\tau(n))$ and hence $o(M_{\mathrm{greedy}}(n))$.
\end{proof}

\begin{remark}
The key geometric input is $\omega(G_X)\le n+1$, which contrasts sharply with $\tau(n)\sim 2^{0.401n}$.  The logarithmic ratio $\ln\omega/\ln\Delta\sim (\log n)/(0.401n)\to 0$ is what drives the improvement.  The de Bruijn--Erd\H{o}s step (Lemma~\ref{lem:compactness}) is essential: the Bonamy--Kelly--Nelson--Postle theorem is a finite-graph result, and the transfer to arbitrary (possibly infinite) packings requires the compactness principle.
\end{remark}

\subsection*{Section 7b: Reed's Bound — A Sharper Estimate for Moderate Dimensions}

The BKNP bound in Theorem~\ref{thm:upper-bound} is asymptotically the strongest known result, but for moderate dimensions its large absolute constant ($C_0 = 72$) makes it weaker than a classical alternative.  We record this alternative, which gives a cleaner closed-form bound superior to BKNP in every dimension $n$ arising in practice.

\begin{theorem}[Reed's Bound for Sphere Packings]\label{thm:reed}
For all sufficiently large $n$,
\[
  M(n) \;\le\; \left\lceil\frac{\tau(n)+n+2}{2}\right\rceil.
\]
\end{theorem}

\begin{proof}
By Lemma~\ref{lem:compactness}, it suffices to bound $\chi(G_X)$ for every finite packing tangency graph $G_X$.  By Lemma~\ref{lem:graph-bounds}, $G_X$ has $\Delta(G_X)\le\tau(n)$ and $\omega(G_X)\le n+1$.  Reed's theorem~\cite{Reed1998} states that for every graph $G$ with $\Delta(G)\ge\Delta_0(\omega(G))$ (a threshold depending on the clique number),
\[
  \chi(G)\;\le\;\left\lceil\frac{\Delta(G)+\omega(G)+1}{2}\right\rceil.
\]
Since $\tau(n)\ge 2^{0.2075n}$ grows exponentially while $\omega(G_X)\le n+1$ is linear, the condition $\Delta\ge\Delta_0(\omega)$ holds for all sufficiently large $n$.  Substituting $\Delta\le\tau(n)$ and $\omega\le n+1$ gives
\[
  \chi(G_X)\;\le\;\left\lceil\frac{\tau(n)+n+2}{2}\right\rceil.
\]
Taking the supremum over all finite $X$ yields the result.
\end{proof}

\begin{remark}[Comparison of the Two Upper Bounds]\label{rem:comparison}
Theorem~\ref{thm:upper-bound} (BKNP) gives $M(n)\le C\tau(n)\sqrt{\log n/n}$, while Theorem~\ref{thm:reed} gives $M(n)\le\tau(n)/2 + O(n)$.  For large $n$:
\[
  \text{BKNP} \;\sim\; C\tau(n)\sqrt{\frac{\log n}{n}} \;=\; o(\tau(n)),
  \qquad
  \text{Reed} \;\sim\; \frac{\tau(n)}{2}.
\]
BKNP is asymptotically stronger ($o(\tau(n))$ vs.\ $\tau(n)/2$), but the constant $C=72$ means Reed's bound dominates in every dimension accessible by computation.  Concretely, for $n\le 10^5$ and $\tau(n)\le 2^{0.401n}$, Reed gives $M(n)\le\tau(n)/2+O(n)$, whereas BKNP gives $72\tau(n)\sqrt{\log n/n}>\tau(n)/2$ (since $72^2\log n/n > 1/4$ for $n\le 10^5$).  The crossover where BKNP beats Reed occurs around $n\approx 4\ln(144^2)\cdot n/\ln n$, i.e.\ $n$ of order $10^6$ or larger.

The combined best-known upper bound is therefore:
\[
  M(n)\;\le\;\min\!\left(\left\lceil\tfrac{\tau(n)+n+2}{2}\right\rceil,\;C\tau(n)\sqrt{\tfrac{\log n}{n}}\right).
\]
For $n\le 10^5$ the Reed bound dominates; for $n\to\infty$ the BKNP bound dominates.

\textbf{Selected values} using exact kissing numbers where known and $\tau(n)\le 2^{0.401n}$ otherwise:

\smallskip
\begin{center}
\begin{tabular}{r|r|r|r|r}
$n$ & $\tau(n)$ & Greedy $\tau+1$ & Reed $\lceil(\tau+n+2)/2\rceil$ & Improvement factor\\
\hline
$3$ & $12$ & $13$ & $9$ & $1.44\times$\\
$4$ & $24$ & $25$ & $15$ & $1.67\times$\\
$8$ & $240$ & $241$ & $125$ & $1.93\times$\\
$24$ & $196560$ & $196561$ & $98293$ & $\approx 2\times$\\
$n\to\infty$ & $\tau(n)$ & $\tau(n)+1$ & $\tau(n)/2+O(n)$ & $\to 2\times$\\
\end{tabular}
\end{center}
\smallskip

Reed's bound thus cuts the greedy upper bound by a factor approaching $2$, independent of $n$.  The BKNP bound cuts it by a factor $\sqrt{n/\log n}\to\infty$, but only asymptotically.
\end{remark}

\begin{remark}[Barriers to Further Improvement]
Both bounds above ultimately rely on $\tau(n)$, the kissing number, as the main input.  Any improvement of the form $M(n)\le\tau(n)^{1-c}$ for a fixed $c>0$ would require structural properties of contact graphs beyond the degree and clique bounds used here.  Candidate approaches include:
\begin{itemize}
  \item \textbf{Lovász theta / SDP}: The Lov\'asz $\vartheta$-function of the contact graph is related to the Cohn--Elkies LP bound for sphere packings, and might yield $M(n)\le\tau(n)/n^{\varepsilon}$ for some $\varepsilon>0$.
  \item \textbf{Entropy compression in neighborhoods}: Since each neighborhood $N(v)$ is itself a kissing configuration (a unit sphere packing on $S^{n-1}$), the local chromatic structure is recursive and may admit a stronger coloring argument.
  \item \textbf{Improving $\tau(n)$}: Any improvement to the Kabatiansky--Levenshtein kissing number bound would immediately improve both $M(n)$ bounds above.  No improvement has been found since 1978.
\end{itemize}
For the \emph{lower} bound, the barrier theorem (Section~6) shows the existing spectral framework is stuck at surplus $O(n^{2/3})$; beating this requires a new geometric construction.  Consequently, the gap between the polynomial lower bound $M(n)\ge n+\Omega(n^{2/3})$ and the exponential upper bound $M(n)\le O(\tau(n))$ remains the central open problem.
\end{remark}

\subsection*{Section 8: Further Upper Bounds and Open Problems}

We record two additional upper bounds and discuss the principal obstruction to further improvement.

\medskip
\textbf{Notation.}  Write $\tau^+(n)$ for the \emph{one-sided kissing number}: the maximum number of non-overlapping unit spheres whose centres lie in a fixed closed half-space of $\mathbb{R}^n$ and all touch a fixed central unit sphere.  Known values: $\tau^+(1)=1$, $\tau^+(2)=4$, $\tau^+(3)=9$ (see~\cite{BvM2022}).  In general, $\tau^+(n)<\tau(n)$ and $\tau^+(n)/\tau(n)\to 1$ as $n\to\infty$; for small $n$, $\tau^+(n)/\tau(n) = 2/3, 3/4, \ldots$, suggesting $\tau^+(n) \approx \tau(n)\cdot(1-1/(n+1))$.

\begin{theorem}[Brooks Bound and Degeneracy Bound]\label{thm:further}
\begin{enumerate}
\item[\textup{(i)}] For all $n\ge 2$: $M(n)\le \tau(n)$.
\item[\textup{(ii)}] For all $n\ge 1$: $M(n)\le \tau^+(n)+1$.
\end{enumerate}
\end{theorem}

\begin{proof}
Both bounds use Lemma~\ref{lem:compactness} (compactness) and Lemma~\ref{lem:graph-bounds} (geometric parameters).

\textbf{(i).}  By Brooks' theorem~\cite{Brooks1941}, a connected graph $G$ satisfies $\chi(G)\le \Delta(G)$ unless $G$ is a complete graph $K_{\Delta+1}$ or an odd cycle.  The tangency graph $G_X$ has $\Delta(G_X)\le \tau(n)$.  For $n\ge 2$, any complete subgraph $K_m \subseteq G_X$ requires $m$ mutually tangent sphere centres, which form an equidistant set in $\mathbb{R}^n$ of size $m$; since the maximum equidistant set in $\mathbb{R}^n$ has size $n+1$, we have $\omega(G_X)\le n+1$.  For $n\ge 2$, $n+1 < \tau(n)+1$ (since $\tau(2)=6>3$, $\tau(3)=12>4$, etc.), so $G_X$ is not a complete graph $K_{\tau(n)+1}$.  Also $G_X$ is not an odd cycle, since for $n\ge 2$ any non-trivial sphere packing has vertices of degree $\ge 1$; in any connected component that is not a single edge (and single edges have $\chi=2\le\tau(n)$), the maximum degree exceeds $2$ and the graph is not a cycle.  Therefore $\chi(G_X)\le\tau(n)$, giving $M(n)\le\tau(n)$.

\textbf{(ii).}  Fix a unit vector $u\in\mathbb{R}^n$ generic (avoiding any direction parallel to an edge of $G_X$).  Order the vertices of $G_X$ by the linear functional $\ell_u(v)=\langle v,u\rangle$; this defines a total order $v_1<v_2<\cdots<v_N$ on the packing centres.

For each $v_i$, the \emph{back-neighbours} are $B(v_i):=\{v_j: j<i,\{v_i,v_j\}\in E(G_X)\}$, i.e.\ the neighbours of $v_i$ with a smaller $\ell_u$-value.  These back-neighbours are sphere centres at distance $1$ from $v_i$ with $\langle w-v_i,u\rangle<0$, lying in the open half of the unit sphere $S^{n-1}(v_i)$ facing $-u$.  This is exactly a one-sided kissing configuration around $v_i$ in the half-space $\{w:\langle w-v_i,u\rangle\le 0\}$, so $|B(v_i)|\le\tau^+(n)$.

Greedy colouring in the order $v_1,v_2,\ldots$ uses at most $\max_i|B(v_i)|+1\le\tau^+(n)+1$ colours.  By Lemma~\ref{lem:compactness}, $M(n)\le\tau^+(n)+1$.
\end{proof}

\begin{remark}
Bound (i) (Brooks) improves greedy by exactly $1$: $M(n)\le\tau(n)$ vs.\ $M(n)\le\tau(n)+1$.  For $n=3$: Brooks gives $M(3)\le 12$, Reed gives $M(3)\le 9$, so Reed dominates.

Bound (ii) (degeneracy via half-space kissing) is a new geometric bound.  For small $n$ where $\tau^+(n)$ is known:
\smallskip
\begin{center}
\begin{tabular}{r|r|r|r|r|r}
$n$ & $\tau(n)$ & $\tau^+(n)$ & Greedy & \textup{Bound (ii)} & Reed \\
\hline
$2$ & $6$ & $4$ & $7$ & $5$ & $5$ \\
$3$ & $12$ & $9$ & $13$ & $10$ & $9$ \\
\end{tabular}
\end{center}
\smallskip
For $n=2$: bound (ii) matches Reed.  For $n=3$: Reed is slightly tighter.  Asymptotically, since $\tau^+(n)<\tau(n)$ but $\tau^+(n)/\tau(n)\to 1$, bound (ii) is asymptotically dominated by Reed's $\tau(n)/2$.

The interest in bound (ii) is geometric rather than numerical: it expresses the chromatic upper bound purely in terms of the one-sided kissing number, a combinatorial-geometric invariant of $\mathbb{R}^n$ that may be tractable via SDP / spherical code methods independently of the full kissing number.
\end{remark}

\medskip
\textbf{Summary of all upper bounds.}  The complete hierarchy for $M(n)$:
\[
  M(n)\;\le\;\min\bigl(\tau^+(n)+1,\;\tau(n),\;\lceil\tfrac{\tau(n)+n+2}{2}\rceil,\;72\,\tau(n)\sqrt{\tfrac{\log n}{n}}\bigr).
\]
For known small dimensions ($n=2,3$): the minimum is achieved by Reed for $n=3$ and by bound (ii) for $n=2$; for $n\to\infty$, BKNP dominates.

\medskip
\textbf{Open problems.}
\begin{enumerate}
\item[\textup{(1)}] \emph{Polynomial factor improvement}: Does $M(n)\le\tau(n)/n^{\varepsilon}$ hold for some $\varepsilon>0$ and all large $n$?  The BKNP bound gives $M(n)\le\tau(n)\cdot O(\sqrt{\log n/n})$ (which is $\tau(n)/n^{1/2-o(1)}$ if $\log$ is base 2), so the answer is yes with $\varepsilon = 1/2 - o(1)$ from BKNP.  Can one do better, e.g.\ $\varepsilon$ close to $1$?

\item[\textup{(2)}] \emph{Contact-graph structure beyond $(\Delta,\omega)$}: Every contact graph is a \emph{1-distance graph} (edges at distance exactly 1, non-edges at distance $>1$ in $\mathbb{R}^n$).  This geometric constraint is not captured by $\Delta$ and $\omega$ alone.  In particular, the Gram matrix of the edge vectors is positive semidefinite with a specific rank-$n$ structure.  Can this algebraic structure be used in a coloring argument?

\item[\textup{(3)}] \emph{One-sided kissing number}: Determine $\tau^+(n)$ for $n\ge 4$.  If $\tau^+(n)\le(1-c)\tau(n)$ for some absolute constant $c>0$, then bound (ii) gives $M(n)\le(1-c)\tau(n)+1$, a constant-factor improvement independent of $n$.  Current evidence ($\tau^+(2)/\tau(2)=2/3$, $\tau^+(3)/\tau(3)=3/4$) is consistent with $\tau^+(n)/\tau(n)\to 1$, in which case bound (ii) is asymptotically equivalent to the greedy bound.

\item[\textup{(4)}] \emph{Lower bound improvement}: Close the gap $n+\Omega(n^{2/3})\le M(n)\le 2^{0.401n(1+o(1))}$.  The barrier theorem (Section~6) rules out a spectral improvement for the four classical families; a new geometric construction with chromatic surplus $\omega(n^\alpha)$ for $\alpha>2/3$ would be needed.
\end{enumerate}


\begin{corollary}[Main Result]\label{cor:main}
The supremum $M(n)$ of chromatic numbers of unit sphere packings in $\mathbb{R}^n$ satisfies:
\begin{enumerate}
\item[\textup{(i)}] \textbf{Lower bound:}
\[
  \limsup_{n\to\infty}\;\frac{M(n)-n}{n^{2/3}}\;\ge\;1.
\]
In particular, $(M(n)-n)/n^{2/3}$ is at least $(1-1/q+1/q^2)^{1/3}\to 1$ along $n=d_q=q^3-q^2+q$ as $q\to\infty$ through prime powers.
\item[\textup{(ii)}] \textbf{Upper bound:}
\[
  M(n)\;=\;O\!\left(\tau(n)\cdot\sqrt{\frac{\log n}{n}}\right)
  \;=\;2^{0.401n(1+o(1))}\cdot O\!\left(\sqrt{\frac{\log n}{n}}\right).
\]
\item[\textup{(iii)}] \textbf{Barrier:} By the Barrier Theorem~\textup{(Theorem~\ref{thm:barrier})}, none of the four named classical spectral families from finite geometry (Kneser, triangular, symplectic polar, hexagon complement) achieves a positive surplus $H(G)-d>0$ within the Hoffman framework.
\end{enumerate}
\end{corollary}

\begin{proof}
\textbf{Lower bound chain.}
By the Spectral Realization Theorem (Theorem~\ref{thm:spectral-realization}) and the explicit spectral computation in Theorem~\ref{thm:gq-family}, the complement $G_q$ of the collinearity graph of $\mathrm{GQ}(q,q^2)$ embeds as the tangency graph of a unit sphere packing in $\mathbb{R}^{d_q}$, so
\[
M(d_q)\;\ge\;\chi(G_q).
\]
By the Hoffman Bound (Theorem~\ref{thm:hoffman}) and the eigenvalues $k_q=q^4$, $s_q=-q$ of $G_q$,
\[
\chi(G_q)\;\ge\;1-\frac{k_q}{s_q}=q^3+1.
\]
Hence $M(d_q)\ge q^3+1 = d_q+(q^2-q+1)$ (Theorem~\ref{thm:main-packing}).

\textbf{Limsup.}
By Theorem~\ref{thm:limsup}(2), taking $q\to\infty$ through prime powers,
\[
\frac{M(d_q)-d_q}{d_q^{2/3}}\;\ge\;\left(1-\frac{1}{q}+\frac{1}{q^2}\right)^{1/3}\;\to\;1,
\]
which gives $\limsup_{n\to\infty}(M(n)-n)/n^{2/3}\ge 1$.

\textbf{Upper bound.}  This is Theorem~\ref{thm:upper-bound}.

\textbf{Barrier.}  This is Theorem~\ref{thm:barrier}.
\end{proof}

\begin{thebibliography}{9}
\bibitem{Chen2015}
H.~Chen,
\textit{Ball packings with high chromatic numbers from strongly regular graphs},
Discrete Mathematics \textbf{340} (2017), no.~7.
arXiv:1502.02070.

\bibitem{BvM2022}
A.~E.~Brouwer and H.~Van Maldeghem,
\textit{Strongly Regular Graphs},
Cambridge University Press, 2022.

\bibitem{Hoffman1970}
A.~J.~Hoffman,
\textit{On eigenvalues and colorings of graphs},
in Graph Theory and its Applications (B.~Harris, ed.), Academic Press, 1970, pp.\ 79--91.

\bibitem{KL1978}
G.~A.~Kabatiansky and V.~I.~Levenshtein,
\textit{Bounds for packings on the sphere and in space},
Problemy Peredachi Informatsii \textbf{14} (1978), 3--25.

\bibitem{PT1984}
S.~E.~Payne and J.~A.~Thas,
\textit{Finite Generalized Quadrangles},
Research Notes in Mathematics \textbf{110}, Pitman, Boston, 1984.
Second edition: EMS Series of Lectures in Mathematics, European Mathematical Society, Z\"{u}rich, 2009.

\bibitem{dBE1951}
N.~G.~de Bruijn and P.~Erd\H{o}s,
\textit{A colour problem for infinite graphs and a problem in the theory of relations},
Indagationes Mathematicae \textbf{13} (1951), 369--373.

\bibitem{BKNP2022}
M.~Bonamy, T.~Kelly, P.~Nelson, and L.~Postle,
\textit{Bounding $\chi$ by a fraction of $\Delta$ for graphs without large cliques},
Journal of the European Mathematical Society \textbf{26} (2024), no.~5, 1451--1499.
arXiv:1803.01051.
\bibitem{Reed1998}
B.~Reed,
\textit{$\omega$, $\Delta$, and $\chi$},
Journal of Combinatorial Theory, Series~B \textbf{74} (1998), 241--260.

\bibitem{Brooks1941}
R.~L.~Brooks,
\textit{On colouring the nodes of a network},
Proceedings of the Cambridge Philosophical Society \textbf{37} (1941), 194--197.

\end{thebibliography}

\end{document}
